\newpage
\Article{73}
Was wir in den Art.~69---72 dargelegt haben, ist eigentlich unabhängig von der Voraussetzung, dass $m$ eine primäre Zahl ist. Von nun an nehmen wir jedoch wenigstens an, dass $a$ ungerade und $b$ gerade ist und außerdem $a,b$ und $a-b$ positive Zahlen sind. Zunächst müssen wir die Grenzen der Werte von $x$ in der Sammlung $G$ bestimmen.

\SourceLeaf{154}%
Setzt man $ay-bx=\eta,(a+b)x-(a-b)y=\zeta,p-2ax-2by=\theta$, so besteht das Kriterium für die Zahlen $x+yi$, die zur Sammlung $G$ gehören, in den drei Bedingungen, dass $\eta,\zeta,\theta$ positive Zahlen sind. Da $px=(a-b)\eta+a\zeta,p(a-2x)=a\theta+2b\eta$ gilt, müssen offenbar $x$ und $2a-x$ positive Zahlen sein, oder $x$ muss einer der Zahlen $1,2,3\ldots\frac12(a-1)$ gleich sein. Da ferner $(a-b)\theta=2b\zeta+p(a-b-2x)$ gilt, ist klar: Solange $x$ kleiner als $\frac12(a-b)$ ist, schließt die zweite Bedingung (nach der $\zeta$ positiv sein muss) die dritte (dass $\theta$ positiv sein muss) bereits ein. Ist dagegen $x$ größer als $\frac12(a-b)$, so ist die zweite Bedingung bereits in der dritten enthalten. Für die Werte von $x$, nämlich $1,2,3\ldots\frac12(a-b-1)$, braucht man daher nur dafür zu sorgen, dass $\eta$ und $\zeta$ positiv werden, oder dass $y$ größer als $\frac{bx}{a}$ und kleiner als $\frac{(a+b)x}{a-b}$ ist. Für einen solchen gegebenen Wert von $x$ beträgt die Gesamtzahl der Zahlen $x+yi$ somit
\[
\left[\frac{(a+b)x}{a-b}\right]-\left[\frac{bx}{a}\right]
\]
wenn wir die Klammern in derselben Bedeutung verwenden, in der wir sie bereits an verschiedenen anderen Stellen gebraucht haben (vgl. \emph{Theorematis arithm. dem. nova}, Art.~4, und \emph{Theorematis fund. in doctr. de residuis quadr.} usw., \emph{Algorithm. nov.}, Art.~3). Für die Werte von $x$, nämlich $\frac12(a-b+1),\frac12(a-b+3)\dotsfour\frac12(a-1)$, genügt es dagegen, $\eta$ und $\theta$ positive Werte zu geben, oder $y$ größer als $\frac{bx}{a}$ und kleiner als $\frac{p-2ax}{2b}$ oder $\frac12b+\frac{aa-2ax}{2b}$ zu machen. Für einen solchen gegebenen Wert von $x$ beträgt die Gesamtzahl der Zahlen $x+yi$ daher
\[
\left[\frac12b+\frac{aa-2ax}{2b}\right]-\left[\frac{bx}{a}\right]
\]
Daraus folgt, dass die Anzahl der Zahlen der Sammlung $G$
\[
g=\Sigma\left[\frac{(a+b)x}{a-b}\right]+\Sigma\left[\frac12b+\frac{aa-2ax}{2b}\right]-\Sigma\left[\frac{bx}{a}\right]
\]
ist; dabei erstreckt sich die erste Summe über alle ganzzahligen Werte von $x$ von $1$ bis $\frac12(a-b-1)$, die zweite von $\frac12(a-b+1)$ bis $\frac12(a-1)$ und die dritte von $1$ bis $\frac12(a-1)$.

\Needspace{14\baselineskip}
Verwenden wir das Zeichen $\varphi$ in derselben Bedeutung wie an der angeführten Stelle (\emph{Theorematis fund.} usw., \emph{Algor. nov.}, Art.~3), sodass also
\[
\varphi(t,u)=\left[\frac{u}{t}\right]+\left[\frac{2u}{t}\right]+\left[\frac{3u}{t}\right]\dotsfour+\left[\frac{t'u}{t}\right]
\]
gilt, wobei $t,u$ beliebige positive Zahlen und $t'$ die Zahl $[\frac12t]$ bezeichnen, so wird jener erste Term $=\varphi(a-b,a+b)$, der dritte $=-\varphi(a,b)$; der zweite wird dagegen
\SourceLeaf{155}%
\[
=\frac14bb+\Sigma\left[\frac{aa-2ax}{2b}\right]
\]
Schreibt man jedoch die Terme in umgekehrter Reihenfolge, so gilt
\[
\begin{aligned}
\Sigma\left[\frac{aa-2ax}{2b}\right]&=\left[\frac{a}{2b}\right]+\left[\frac{3a}{2b}\right]+\left[\frac{5a}{5b}\right]+\dotsfour+\left[\frac{(b-1)a}{2b}\right]\\
&=\varphi(2b,a)-\varphi(b,a)
\end{aligned}
\]
Unsere Formel nimmt daher die folgende Gestalt an:
\[
g=\varphi(a-b,a+b)+\varphi(2b,a)-\varphi(a,b)-\varphi(b,a)+\frac14bb
\]

Betrachten wir zunächst den Term $\varphi(a-b,a+b)$, der unmittelbar in $\varphi(a-b,2b)+1+2+3+\text{ usw. }+\frac12(a-b-1)$ oder in
\[
\varphi(a-b,2b)+\frac18((a-b)^2-1)
\]
übergeht.

Da ferner nach dem allgemeinen Satz $\varphi(t,u)+\varphi(u,t)=[\frac12t]\pdot[\frac12u]$ gilt, wenn $t,u$ teilerfremde positive ganze Zahlen sind, haben wir
\[
\varphi(a-b,2b)=\frac12b(a-b-1)-\varphi(2b,a-b)
\]
und somit
\[
\varphi(a-b,a+b)=\frac18(aa+2ab-3bb-4b-1)-\varphi(2b,a-b)
\]

Ordnen wir die Teile von $\varphi(2b,a-b)$ folgendermaßen an:
\[
\begin{aligned}
&\left[\frac{a-b}{2b}\right]+\left[\frac{3(a-b)}{2b}\right]+\left[\frac{5(a-b)}{2b}\right]+\text{ usw. }+\left[\frac{(b-1)(a-b)}{2b}\right]\\
&+\left[\frac{a-b}{b}\right]+\left[\frac{2(a-b)}{b}\right]+\left[\frac{3(a-b)}{b}\right]+\text{ usw. }+\left[\frac{\frac12b(a-b)}{b}\right]
\end{aligned}
\]
Die zweite Reihe wird offenbar
\[
\begin{aligned}
&=\varphi(b,a-b)=\varphi(b,a)-1-2-3-\text{ usw. }-\frac12b\\
&=\varphi(b,a)-\frac18(bb+2b)
\end{aligned}
\]
Die erste Reihe stellen wir nach Umkehrung ihrer Gliederfolge so dar:
\[
\begin{aligned}
&\left[\frac12(a+1-b)-\frac{a}{2b}\right]+\left[\frac12(a+3-b)-\frac{3a}{2b}\right]\\
&+\left[\frac12(a+5-b)-\frac{5a}{2b}\right]+\text{usw.}+\left[\frac12(a-1)-\frac{(b-1)a}{2b}\right]
\end{aligned}
\]
Dieser Ausdruck. Da für eine ganze Zahl $t$ und eine gebrochene Zahl $u$ allgemein $[t-u]=t-1-[u]$ gilt, geht er in den folgenden über:
\[
\begin{aligned}
&\frac18b(2a-4-b)-\left[\frac{a}{2b}\right]-\left[\frac{3a}{2b}\right]-\left[\frac{5a}{2b}\right]-\text{ usw. }-\left[\frac{(b-1)a}{2b}\right]\\
&=\frac18b(2a-4-b)-\varphi(2b,a)+\varphi(b,a)
\end{aligned}
\]

\SourceLeaf{156}%
Daraus folgt
\[
\varphi(2b,a-b)=2\varphi(b,a)-\varphi(2b,a)+\frac14b(a-3-b)
\]
und daher
\[
\varphi(a-b,a+b)=\varphi(2b,a)-2\varphi(b,a)+\frac18(aa-bb+2b-1)
\]
Setzen wir diesen Wert in die oben angegebene Formel für $g$ ein und außerdem $\varphi(a,b)+\varphi(b,a)=\frac14b(a-1)$, so erhalten wir
\[
g=2\varphi(2b,a)-2\varphi(b,a)+\frac18(aa-2ab+bb+4b-1)
\]

\newpage
\Article{74}
Durch genau dieselben Überlegungen erledigt sich der Fall, in dem $a$, $b$ positiv bleiben, aber $a-b$ negativ oder $b-a$ positiv ist. Die Gleichungen $p(a-2x)=2b\eta+a\theta$, $p(b-a+2x)=2b\zeta+(b-a)\theta$ zeigen, dass $\frac12a-x$ und $x+\frac12(b-a)$ positiv sein müssen und somit $x$ einer der Zahlen $-\frac12(b-a-1),-\frac12(b-a-3),-\frac12(b-a-5)\dotsfour+\frac12(a-1)$ gleich sein muss. Ferner folgt aus der Gleichung $px+(b-a)\eta=a\zeta$, dass für negative Werte von $x$ die Bedingung, nach der $\eta$ positiv sein muss, bereits in der Bedingung enthalten ist, nach der $\zeta$ positiv sein muss; das Umgekehrte gilt, wenn $x$ ein positiver Wert gegeben wird. Die Werte von $y$ müssen daher für einen bestimmten negativen Wert von $x$ zwischen $\frac{(a+b)x}{a-b}$ und $\frac{p-2ax}{2b}$, für einen positiven Wert von $x$ dagegen zwischen $\frac{bx}{a}$ und $\frac{p-2ax}{2b}$ liegen. Für $x=0$ sind diese Grenzen offenbar $0$ und $\frac{p-2ax}{2b}$, wobei der Wert $y=0$ selbst ausgeschlossen ist. Daraus folgt
\[
g=-\Sigma\left[\frac{(a+b)x}{a-b}\right]+\Sigma\left[\frac12b+\frac{aa-2ax}{2b}\right]-\Sigma\left[\frac{bx}{a}\right]
\]
wobei sich die erste Summe über alle negativen Werte von. $x$ von $-1$ bis $-\frac12(b-a-1)$ erstreckt, die zweite über alle Werte von $x$ von $-\frac12(b-a-1)$ bis $\frac12(a-1)$ und die dritte über alle positiven Werte von $x$ von $+1$ bis $\frac12(a-1)$. So ergibt die erste Summe $-\varphi(b-a,b+a)$, die zweite wie im vorigen Artikel $\frac14bb+\varphi(2b,a)-\varphi(b,a)$ und schließlich die dritte $-\varphi(a,b)$, also gilt
\[
g=-\varphi(b-a,b+a)+\varphi(2b,a)-\varphi(b,a)-\varphi(a,b)+\frac14bb
\]
Nun erhält man auf dieselbe Weise wie im vorigen Artikel
\SourceLeaf{157}%
\[
\begin{aligned}
\varphi(b-a,b+a)&=\varphi(b-a,2b)-\frac18((b-a)^2-1)\\
&=\frac18(3bb-2ab-aa-4b+1)-\varphi(2b,b-a)
\end{aligned}
\]
sowie
\[
\varphi(2b,b-a)=\varphi(2b,a)-2\varphi(b,a)+\frac14b(b-1-a)
\]
und somit
\[
\varphi(b-a,b+a)=2\varphi(b,a)-\varphi(2b,a)+\frac18(bb-aa-2b+1)
\]
und schließlich
\[
g=2\varphi(2b,a)-2\varphi(b,a)+\frac18(aa-2ab+bb+4b-1)
\]
Damit ist bewiesen, dass dieselbe Formel für $g$ gilt, ob $a-b$ positiv oder negativ ist, sofern $a$, $b$ positiv sind.

\Article{75}
Um eine weitere Reduktion zu erreichen, setzen wir
\[
\begin{aligned}
L&=\left[\frac{a}{2b}\right]+\left[\frac{2a}{2b}\right]+\left[\frac{3a}{2b}\right]+\text{usw.}+\left[\frac{\frac12ba}{2b}\right]\\
M&=\left[\frac{(\frac12b+1)a}{2b}\right]+\left[\frac{(\frac12b+2)a}{2b}\right]+\left[\frac{(\frac12b+3)a}{2b}\right]+\text{usw.}+\left[\frac{ba}{2b}\right]\\
N&=\left[\frac{a+b}{2b}\right]+\left[\frac{2a+b}{2b}\right]+\left[\frac{3a+b}{2b}\right]+\text{usw.}+\left[\frac{\frac12ba+b}{2b}\right]
\end{aligned}
\]
Man sieht leicht, dass allgemein $[u]+[u+\frac12]=[2u]$ gilt, welche reelle Größe $u$ auch bezeichnen mag. Daher gilt $L+N=\varphi(b,a)$, und da offenbar $L+M=\varphi(2b,a)$ ist, folgt
\[
\varphi(2b,a)-\varphi(b,a)=M-N
\]

Ferner ist ersichtlich, dass die Summe aus dem ersten Glied der Reihe $N$ und dem vorletzten Glied der Reihe $M$, nämlich $\left[\frac{a+b}{2b}\right]+\left[\frac{(b-1)a}{2b}\right]$, $=\frac12(a-1)$ wird und dass dieselbe Summe aus dem zweiten Glied der Reihe $N$ und dem drittletzten Glied der Reihe $M$ entsteht, und so weiter. Da auch das letzte Glied der Reihe $M$ $=\frac12(a-1)$ ist, das letzte Glied der Reihe $N$ dagegen $=\left[\frac{a+2}{4}\right]=\frac14(a\mp1)$, wobei das obere oder untere Vorzeichen gilt, je nachdem $a$ die Form $4n+1$ oder $4n-1$ hat, erhalten wir
\[
M+N=\frac14(a-1)b+\frac14(a\mp1)
\]
und daher
\[
\varphi(2b,a)-\varphi(b,a)=\frac14(a-1)b+\frac14(a\mp1)-2N
\]

\SourceLeaf{158}%
Die in den Art.~73 und 74 gefundene Formel für $g$ geht somit in die folgende über:
\[
g=\frac18((a+b)^2-1)+2n-4N
\]
wobei $a\mp1=4n$ gesetzt wird und $n$ eine ganze Zahl ist. Da daraus $1=16nn-8an+aa$ folgt, lässt sich diese Formel auch so darstellen:
\[
g=\frac18(-aa+2ab+bb+1)+4(\frac12(a+1)n-nn-N)
\]
Da $g$ der Charakter der Zahl $-1-i$ für den Modul $a+bi$ ist, ist dieser Charakter daher $\equiv\frac18(-aa+2ab+bb+1)\pmod4$. Das ist genau der oben (Art.~64) durch Untersuchung von Beispielen gefundene Satz; daraus folgen unmittelbar die Sätze über die Charaktere der Zahlen $1+i$, $1-i$, $-1+i$. Somit sind diese vier Sätze für den Fall, dass $a$ und $b$ positiv sind, nun streng bewiesen.

\Needspace{16\baselineskip}
\Article{76}
Bleibt $a$ positiv und ist $b$ negativ, so setzen wir $b=-b'$, sodass $b'$ positiv wird. Es ist nun bewiesen, dass für den Modul $a+b'i$ der Charakter der Zahl $-1-i$ $\equiv\frac18(-aa+2ab'+b'b'+1)\pmod4$ ist. Daher ist der Charakter der Zahl $-1+i$ für den Modul $a-b'i$ nach dem in Art.~62 angegebenen Satz $\equiv\frac18(aa-2ab'-b'b'-1)$, das heißt, der Charakter der Zahl $-1+i$ für den Modul $a+bi$ wird $\equiv\frac18(aa+2ab-bb-1)$. Das ist genau der in Art.~64 angegebene Satz, aus dem die drei übrigen über die Charaktere der Zahlen $1+i$, $1-i$, $-1-i$ unmittelbar folgen. Somit sind diese Sätze auch für den Fall bewiesen, dass $b$ negativ ist, also für alle Fälle, in denen $a$ positiv ist.

Ist schließlich $a$ negativ, so setzen wir $a=-a'$, $b=-b'$. Nach dem bereits Bewiesenen ist der Charakter der Zahl $1+i$ bezüglich des Moduls $a'+b'i$ $\equiv\frac18(-a'a'+2a'b'-3b'b'+1)\pmod4$, und es ist gleichgültig, ob wir die Zahl $a'+b'i$ oder ihre entgegengesetzte Zahl $-a'-b'i$ als Modul nehmen. Also ist der Charakter der Zahl $1+i$ bezüglich des Moduls $a+bi$ offenbar $\equiv\frac18(-aa+2ab-3bb+1)$, und Entsprechendes gilt für die Charaktere der Zahlen $1-i$, $-1+i$, $-1-i$.

Daraus schließen wir, dass der Beweis der Sätze über die Charaktere der Zahlen $1+i$, $1-i$, $-1+i$, $-1-i$ (Art.~63. 64) keiner Einschränkung mehr unterliegt.
\BookRule{.28\linewidth}
